% -------------------------------------------------------------
%    Conclusion
% -------------------------------------------------------------

\section{Conclusion}

In the process of developing an application that allows the user to optimize the time schedule of heat producing, as well as electricity requiring and producing units, for cost-effective. 

The key points of the system were Design and implementation of:

\begin{itemize}
    \item \texttt{Data flow}, deserializing and/or passing Data from the application to the \texttt{Models}, along with visualizing them.

    \item Implementation of logic and calculation of the \texttt{optimizing process}

    \item Compiling all functionalities in \texttt{one Window}, while keeping it user-friendly and non-complex 
\end{itemize}

By calculating a value for each production unit, for their own cost per hour to produce heat, based on the net production cost of each unit depending on heat demand and electricity price per hour, and sorting a list of available units for that value, the application can create cost-effective time schedules. 
Additionally, the calculation considers settings such as: Scenarios, periods and maintenance, which can easily be changed on the main window. 

For future development, the application could receive a change for more flexibility, especially for the visualisation of the units, since the current version works inefficient for a higher number of units.
Another future approach would be the implementation of an API to improve the Data flow. 
Additionally, the application could be improved to operate for different cities. 
